% Estilo de los reportes del proyecto, con la identidad visual de la Universidad
% de los Andes: amarillo, negro y blanco, y tipografía sin serifas (ver README.md
% de esta carpeta). Lo comparten los reportes de las dos entregas.
%
% El reporte define antes: \titulo, \titulocorto, \subtitulo, \entrega, \fecha,
% \urlrepositorio, \urlvideo y \maxpaginas, y pone \input{../../identidad/estilo-reporte}.
% Se compila con LuaLaTeX.

\usepackage[spanish,es-tabla,es-nodecimaldot]{babel}
\usepackage{fontspec}
\usepackage[letterpaper,left=2.5cm,right=2.5cm,top=2.6cm,bottom=2.4cm,
  headheight=14pt,headsep=0.8cm,footskip=1.0cm]{geometry}
\usepackage{graphicx}
\usepackage{booktabs}
\usepackage{enumitem}
\usepackage{titlesec}
\usepackage{csquotes}
\usepackage{xcolor}
\usepackage{luacolor}
\usepackage{lua-ul}
\usepackage{tikz}
\usepackage{fancyhdr}
\usepackage{lastpage}
\usepackage[font=small,labelfont=bf,labelsep=period]{caption}
\usepackage[backend=biber,style=numeric,sorting=none]{biblatex}
\usepackage[hidelinks]{hyperref}

\newcommand{\dirIdentidad}{../../identidad}

% --- Colores -----------------------------------------------------------------
% Los oficiales del logosímbolo: negro y amarillo (C0 M0 Y100 K0). El amarillo se
% define en CMYK para que sea exactamente el del logo.
\definecolor{amarillo}{cmyk}{0,0,1,0}
\colorlet{negro}{black}
\definecolor{gris}{gray}{0.40}
\definecolor{grisclaro}{gray}{0.94}

% --- Tipografía --------------------------------------------------------------
% Arial 12 (lo exige el enunciado; el manual recomienda fuentes sin serifas). Si
% Arial no está instalada se usa una fuente con las mismas medidas, para que el
% número de páginas no cambie, pero el PDF que se entrega debe salir con Arial.
\IfFontExistsTF{Arial}{\setmainfont{Arial}}{%
  \IfFontExistsTF{Liberation Sans}{\setmainfont{Liberation Sans}}{\setmainfont{TeX Gyre Heros}}%
  \AtEndDocument{\PackageWarningNoLine{reporte}{Arial no esta instalada: este PDF usa una
    fuente sustituta. Compile la version que se entrega en un equipo con Arial}}}
\renewcommand{\familydefault}{\rmdefault}

% --- Títulos, listas y párrafos: compactos, el espacio es para el contenido -----
\titleformat{\section}{\large\bfseries}
  {\colorbox{amarillo}{\makebox[1.3em]{\strut\thesection}}}{0.6em}{}
\titleformat{\subsection}{\normalsize\bfseries}{\thesubsection}{0.5em}{}
\titlespacing*{\section}{0pt}{1.4ex plus 0.4ex}{0.7ex}
\titlespacing*{\subsection}{0pt}{1ex plus 0.3ex}{0.4ex}
\setlist{nosep,leftmargin=*}
\setlist[itemize,1]{label=\textcolor{amarillo}{\rule[0.15ex]{0.9ex}{0.9ex}}}
\setlength{\parskip}{0.5ex}
\setlength{\parindent}{0pt}

% --- Enlaces: texto negro con subrayado amarillo ---------------------------------
\newcommand{\enlace}[2]{\href{#1}{\underLine[color=amarillo,top=-0.22ex,bottom=-0.42ex]{#2}}}
% Enlace a un archivo o carpeta del repositorio: \repo{ruta/al/archivo}{texto}.
\newcommand{\repo}[2]{\enlace{\urlrepositorio/blob/main/#1}{#2}}

% --- Recuadro para el resultado o la cifra que no se puede pasar por alto --------
% Fondo gris con una barra amarilla a la izquierda. No se parte entre páginas.
\newsavebox{\cajadestacado}
\newenvironment{destacado}{%
  \par\addvspace{1.2ex}\noindent\setlength{\fboxsep}{6pt}%
  \begin{lrbox}{\cajadestacado}%
  \begin{minipage}{\dimexpr\linewidth-3pt-2\fboxsep\relax}\setlength{\parskip}{0.5ex}}
  {\end{minipage}\end{lrbox}%
  \textcolor{amarillo}{\vrule width 3pt}\colorbox{grisclaro}{\usebox{\cajadestacado}}%
  \par\addvspace{1.2ex}}

% --- Encabezado y pie de página --------------------------------------------------
% Franja amarilla de borde a borde, con una línea negra debajo.
\newcommand{\franja}{%
  \begin{tikzpicture}[remember picture,overlay]
    \fill[amarillo] (current page.north west) rectangle
      ([yshift=-1.15cm]current page.north east);
    \fill[negro] ([yshift=-1.15cm]current page.north west) rectangle
      ([yshift=-1.21cm]current page.north east);
    \node[anchor=west,font=\small\bfseries,text=negro]
      at ([xshift=2.5cm,yshift=-0.6cm]current page.north west) {\titulocorto};
    \node[anchor=east,font=\small,text=negro]
      at ([xshift=-2.5cm,yshift=-0.6cm]current page.north east) {\entrega};
  \end{tikzpicture}}
\pagestyle{fancy}
\fancyhf{}
\fancyhead[L]{\franja}
\fancyfoot[L]{\footnotesize\color{gris}MINE-4101 Ciencia de Datos Aplicada · Universidad de los Andes}
\fancyfoot[R]{\footnotesize\color{gris}\thepage\ de \pageref*{LastPage}}
\renewcommand{\headrulewidth}{0pt}
\renewcommand{\footrulewidth}{0.4pt}

% --- Portada ----------------------------------------------------------------------
% El logo a color va sobre blanco, sin modificar y acompañado del texto legal, como
% pide el manual de identidad.
\newcommand{\legaluniandes}{Universidad de los Andes | Vigilada Mineducación.
  Reconocimiento como Universidad: Decreto 1297 del 30 de mayo de 1964.
  Reconocimiento personería jurídica: Resolución 28 del 23 de febrero de 1949 Minjusticia.}

\newcommand{\portada}{%
  \begin{titlepage}
  \begin{tikzpicture}[remember picture,overlay]
    \node[anchor=north west,inner sep=0pt]
      at ([xshift=2.5cm,yshift=-2.2cm]current page.north west)
      {\includegraphics[height=1.9cm]{\dirIdentidad/logo-uniandes.pdf}};
    % bloque amarillo de borde a borde, con el título
    \fill[amarillo] ([yshift=-7.6cm]current page.north west) rectangle
      ([yshift=-16.0cm]current page.north east);
    \node[anchor=north west,inner sep=0pt,text width=16.5cm,align=left]
      at ([xshift=2.5cm,yshift=-8.7cm]current page.north west) {%
        \hyphenpenalty=10000\exhyphenpenalty=10000
        {\small\bfseries\addfontfeature{LetterSpace=12}\MakeUppercase{Proyecto final · \entrega}}\\[0.9cm]
        {\fontsize{30}{35}\selectfont\bfseries\titulo\par}\vspace{0.6cm}
        {\Large\subtitulo\par}};
    % franja negra con el curso
    \fill[negro] ([yshift=-16.0cm]current page.north west) rectangle
      ([yshift=-17.3cm]current page.north east);
    \node[anchor=west,text=white,font=\normalsize\bfseries]
      at ([xshift=2.5cm,yshift=-16.65cm]current page.north west)
      {MINE-4101 Ciencia de Datos Aplicada};
    \node[anchor=east,text=white,font=\normalsize]
      at ([xshift=-2.5cm,yshift=-16.65cm]current page.north east) {Semestre 2026-20};
    % integrantes y profesores
    \node[anchor=north west,inner sep=0pt,text width=8.5cm,align=left]
      at ([xshift=2.5cm,yshift=-18.6cm]current page.north west) {%
        {\small\bfseries\addfontfeature{LetterSpace=12}INTEGRANTES}\\[0.25cm]
        Daniel Felipe Triviño Santana\\ Nicolas Bedoya Figueroa\\
        Joubert Daniel Álvarez\\ Gabriel Bejarano};
    \node[anchor=north west,inner sep=0pt,text width=7cm,align=left]
      at ([xshift=12cm,yshift=-18.6cm]current page.north west) {%
        {\small\bfseries\addfontfeature{LetterSpace=12}PROFESORES}\\[0.25cm]
        Fabian Camilo Peña\\ Juan Pablo Reyes\\[0.5cm]
        \enlace{\urlrepositorio}{Repositorio} · \enlace{\urlvideo}{Video}\\[0.15cm]
        \fecha};
    % texto legal que acompaña al logo
    \draw[negro,line width=0.4pt] ([xshift=2.5cm,yshift=2.5cm]current page.south west)
      -- ([xshift=-2.5cm,yshift=2.5cm]current page.south east);
    \node[anchor=north west,inner sep=0pt,text width=16.5cm,align=left,
      font=\scriptsize,text=gris]
      at ([xshift=2.5cm,yshift=2.3cm]current page.south west) {\legaluniandes};
  \end{tikzpicture}
  \end{titlepage}
  \setcounter{page}{2}}

% Título compacto, para cuando no hay portada.
\newcommand{\titulocompacto}{%
  \begin{center}
    {\Large\bfseries\titulo\par}\vspace{0.5ex}
    \entrega\ · MINE-4101 Ciencia de Datos Aplicada · 2026-20\\[0.5ex]
    Daniel Felipe Triviño Santana, Nicolas Bedoya Figueroa,\\
    Joubert Daniel Álvarez, Gabriel Bejarano\\[0.5ex]
    \enlace{\urlrepositorio}{Repositorio} · \enlace{\urlvideo}{Video de sustentación}
  \end{center}}

% \encabezado pone la portada o, con \portadafalse, el título compacto.
\newif\ifportada
\portadatrue
\newcommand{\encabezado}{\ifportada\portada\else\titulocompacto\fi}

% --- Aviso si el reporte se pasa del máximo de páginas ---------------------------
\AtEndDocument{\ifnum\value{page}>\maxpaginas\relax
  \PackageWarningNoLine{reporte}{El reporte tiene \thepage\space paginas y el maximo
    es \maxpaginas}\fi}
